\section{Matter carriers, Gauss constraints, and defect fusion}
\label{sec:matter-textbook}

\subsection{Matter fields and covariant transport}

Chapter~\ref{ch:1} introduced a representation on a finite-dimensional vector space $W$
over a field $k$ as a homomorphism
$\rho:G\to\operatorname{GL}(W)$.  When $k=\mathbb C$ and $W$ carries an inner
product, a unitary representation lands in $U(W)$, the group of linear maps
preserving that inner product.

Finite-dimensional does not mean finite as a set when $k$ is infinite.  The
representation carrier $W$ is parameter data on which linear calculations are
performed.  When a profile is required to have a finite state type, it must
separately declare a finite $G$-stable palette $P\subseteq W$, a finite union
of $G$-orbits in $W$, or another finite $G$-set equipped with an equivariant
map to $W$.  The full vector space may instead be retained, in a profile
that Definition~\ref{def:gauge-profile} calls finite-base: its canvas and
gauge group are finite, while its state type need not be.  A profile whose
state type is finite as well is finite-state.

The intrinsic matter carrier at $v$ is the \emph{associated fibre}
\[
 E_v\defeq T_v\times_G W,
\]
the set of equivalence classes $[t,w]$ under
$(t\cdot g,w)\sim(t,\rho(g)w)$.  An intrinsic site matter field is a choice
$\Psi_v\in E_v$.  After choosing $t_v\in T_v$, write
$\Psi_v=[t_v,\psi_v]$ with coordinate $\psi_v\in W$.  If the reference changes
passively to $t_v\cdot h_v^{-1}$, the same intrinsic vector has coordinate
$\rho(h_v)\psi_v$.  An active torsor automorphism $\alpha_v$ induces the
well-defined associated-fibre map
$[t,w]\mapsto[\alpha_v(t),w]$.  If
$\alpha_v(t_v)=t_v\cdot h_v$, its image in the fixed trivialization also has
coordinate
\begin{equation}
\psi_v^h=\rho(h_v)\psi_v.
 \label{eq:matter-gauge-action}
\end{equation}
The connection induces an intrinsic map
$[t,w]\mapsto[\mathcal U_e(t),w]$ from $E_u$ to $E_v$.  In local coordinates,
transport along $e:u\to v$ therefore sends the source vector to
\[
 \rho(U_e)\psi_u\in W.
\]
This transported vector belongs to the target frame:
\begin{equation}
\rho(U_e^h)\psi_u^h
 =\rho(h_v)\rho(U_e)\psi_u.
 \label{eq:matter-covariance}
\end{equation}
The equality is \emph{covariance}, not invariance.  Both sides change, but
they change in the same target-frame manner.

\begin{example}[The one-link $C_3$ calculation with matter]
Let $\omega=e^{2\pi i/3}$ and use the one-dimensional representation
$\rho(r)=\omega$.  In the coordinates of Example~\ref{ex:c3-two-readings}, take
$U_e=r$ and $\psi_u=1$.  The transported vector is $\omega$.  Under
$h_u=r^2$ and $h_v=r$, the new link coordinate is $U_e^h=1$ and the new
source coordinate is $\psi_u^h=\omega^2$.  Transport now gives $\omega^2$.
This agrees with transforming the old target vector:
$\rho(h_v)\omega=\omega^2$.  The numerical coordinate changed, while the
covariant transport statement did not.
\end{example}

\subsection{Invariant vertex couplings}

Retain the chosen orientation set $E^+(X)$.  From here on the letter $W$ is
kept for edge carriers: the matter carrier at a vertex, written $W$ with
representation $\rho$ in the preceding subsection, is written $M_v$ with
representation $\mu_v$.  Assign a
declared representation
\[
 \rho_e:G\longrightarrow\operatorname{GL}(W_e)
\]
to each $e\in E^+(X)$, and assign a declared matter representation
$\mu_v:G\to\operatorname{GL}(M_v)$ to each vertex.  The represented link
transport is a linear map from the
source copy of $W_e$ to the target copy, and therefore has tensor type
\[
 \rho_e(U_e)\in\Hom(W_e,W_e)\cong W_e\otimes W_e^*.
\]
The first factor is read at the target and the dual factor at the source.
Thus an incoming edge contributes $W_e$ at $v$, while an outgoing edge
contributes the dual representation $W_e^*$, whose action is contragredient.
Explicitly, if $\ell\in W_e^*$, then
$(g\cdot\ell)(w)=\ell(\rho_e(g)^{-1}w)$.
The total carrier seen at $v$ is
\[
 \mathcal K_v=
 \left(\bigotimes_{\substack{e\in E^+(X)\\t(e)=v}}W_e\right)
 \otimes
 \left(\bigotimes_{\substack{e\in E^+(X)\\s(e)=v}}W_e^*\right)
 \otimes M_v.
\]
The group acts diagonally on this tensor product, using $\rho_e$ on incoming
factors, the contragredient representation on outgoing factors, and $\mu_v$
on matter.  For two $G$-representations $A,B$, the notation $\Hom_G(A,B)$
means the vector space of equivariant linear maps $A\to B$.

\begin{definition}[Vertex-invariance witness types]
The \emph{invariant-tensor type} at $v$ is
\begin{equation}
\mathcal I_v\defeq\Hom_G(\mathbf 1,\mathcal K_v)
 \cong \mathcal K_v^G,
 \label{eq:gauss-intertwiner}
\end{equation}
where $\mathbf 1$ is the trivial one-dimensional representation and
$\mathcal K_v^G$ is the subspace fixed by the diagonal $G$-action.
An inhabitant $i_v\in\mathcal I_v$ is an invariant tensor joining the incident
edge and matter carriers.  When a covector is required for contraction, its
type is
\[
 \mathcal I_v^{\vee}\defeq\Hom_G(\mathcal K_v,\mathbf1).
\]
A declared invariant inner product on $\mathcal K_v$ identifies
$\mathcal I_v$ with $\mathcal I_v^{\vee}$; without one the two types stay
distinct.  These tensors are also called
\emph{Gauss intertwiners}; a graph whose edges carry representations and
whose vertices carry such intertwiners is a spin network in the sense of
Baez \cite{BaezSpinNetworks}; below, ``Gauss constraint'' will mean a predicate on
the full state rather than an intertwiner type.
\end{definition}

\begin{example}[$C_2$ parity at one vertex]
\label{ex:c2-vertex-parity}
Let $G=C_2$ and let $\varepsilon$ be its one-dimensional sign
representation.  The tensor power $\varepsilon^{\otimes n}$ is trivial
exactly when $n$ is even.  Thus three incident sign-labeled edges with no
matter give
$\Hom_G(\mathbf1,\varepsilon^{\otimes3})=0$: no invariant vertex witness
exists.  Adding
one sign-labeled matter carrier gives
$\varepsilon^{\otimes4}\cong\mathbf1$, so the witness space is
one-dimensional.  In this example kinematic gauge neutrality is even
parity.
\end{example}

\begin{example}[A non-Abelian invariant contraction]
Let $S_3$ act on
\[
 V_{\rm st}=\{(x_1,x_2,x_3)\in\mathbb C^3:x_1+x_2+x_3=0\}
\]
by permuting coordinates.  This is the two-dimensional standard
representation.  The bilinear map
\[
 \beta:V_{\rm st}\otimes V_{\rm st}\longrightarrow\mathbb C,
 \qquad
 \beta(x,y)=x_1y_1+x_2y_2+x_3y_3
\]
is $S_3$-equivariant when $\mathbb C$ carries the trivial action, because a
permutation merely reorders the three summands.  Hence
$\beta\in\Hom_{S_3}(V_{\rm st}\otimes V_{\rm st},\mathbf1)$.  For example,
$\beta((1,-1,0),(0,1,-1))=-1$, and applying any common permutation to both
vectors leaves that scalar unchanged.  This is an explicit non-Abelian
vertex-invariance witness; it is still not, by itself, a predicate on a full
field state.
\end{example}

A \emph{finite tensor network} here means a finite collection of tensors whose
paired carrier and dual-carrier indices are joined by evaluation maps.  It is
a scalar network when no unpaired index remains.

\begin{proposition}[Invariant matter contraction]
\label{prop:matter-transport}
\status{PROVED-LOCAL}
Let the declared edge and vertex representations be finite-dimensional over
one fixed coefficient field.  Representation-valued edge transport is gauge
covariant.  Any finite scalar
tensor network built from these transport tensors, vertex-invariance witnesses
$i_v\in\mathcal I_v$ or $i_v^{\vee}\in\mathcal I_v^{\vee}$, and $G$-equivariant
evaluation or contraction maps is
gauge invariant.
\end{proposition}

\begin{proof}
Equation~\eqref{eq:matter-covariance} proves covariance on each transported
edge.  At a vertex, all remaining factors of $h_v$ act diagonally on
$\mathcal K_v$.  Both $i_v:\mathbf1\to\mathcal K_v$ and
$i_v^{\vee}:\mathcal K_v\to\mathbf1$ are intertwiners.  Equivariant evaluation
and contraction maps commute with the action, so paired carrier and
dual-carrier factors cancel.  After all indices have been contracted to the trivial
representation, no local gauge factor remains; the resulting scalar is
constant on gauge orbits.
\end{proof}

\subsection{Gauss constraints are predicates on states}

\begin{definition}[Gauss-constraint profile]
Let a state type $\mathcal S$ carry an action of the intrinsic gauge
group $\mathscr G_X$.  A \emph{Gauss-constraint profile} is a family of
predicates
\[
 \mathsf{Gauss}_v:\mathcal S\longrightarrow\mathsf{Prop}
 \qquad(v\in V(X))
\]
together with coherent witness transports
\[
 \gamma_{\alpha,x,v}:\mathsf{Gauss}_v(x)
   \xrightarrow{\sim}\mathsf{Gauss}_v(\alpha\cdot x),
 \qquad
 \gamma_{1,x,v}=\id,
 \qquad
 \gamma_{\alpha\beta,x,v}
 =\gamma_{\alpha,\beta x,v}\circ\gamma_{\beta,x,v}.
\]
Thus gauge invariance means an isomorphism of witness types, not only equality
of truth values.  In particular,
$\mathsf{Gauss}_v(\alpha\cdot x)$ holds exactly when
$\mathsf{Gauss}_v(x)$ holds.  Its admitted state type is
\[
 \mathcal S_{\mathsf G}
 \defeq
 \{x\in\mathcal S:\mathsf{Gauss}_v(x)\text{ for every }v\}.
\]
An inhabitant of $\mathsf{Gauss}_v(x)$ is a \emph{Gauss-state witness}.  For an
update $F$, a family of implications
\[
 \mathsf{Gauss}_v(x)\Longrightarrow\mathsf{Gauss}_v(F(x))
\]
is a separate \emph{Gauss-preservation witness}.

For a relation $R\subseteq\mathcal S\times\mathcal S$, Gauss preservation
means
\[
 x\in\mathcal S_{\mathsf G},\ (x,y)\in R
 \quad\Longrightarrow\quad y\in\mathcal S_{\mathsf G}.
\]
For a Markov kernel $K$ on the finite state type, it means
\[
 x\in\mathcal S_{\mathsf G}
 \quad\Longrightarrow\quad
 \sum_{y\in\mathcal S_{\mathsf G}}K(x,y)=1.
\]
Equivalently, an admitted input assigns zero probability to every
nonadmitted output.  These are the relational and stochastic
Gauss-preservation obligations used below.
\end{definition}

\begin{proposition}[Constraint descent and preservation]
\label{prop:gauss-constraint}
\status{PROVED-LOCAL}
The admitted state type $\mathcal S_{\mathsf G}$ of a Gauss-constraint profile
is stable under gauge transformations.  If a gauge-covariant update $F$
comes with Gauss-preservation witnesses, its restriction to
$\mathcal S_{\mathsf G}$ is gauge covariant and descends to an update on
$\mathcal S_{\mathsf G}/\mathscr G_X$.  An equivariant Gauss-preserving
relation restricts to an equivariant relation on
$\mathcal S_{\mathsf G}$, and an equivariant Gauss-preserving Markov kernel
restricts to a Markov kernel on $\mathcal S_{\mathsf G}$.  Their orbit
descents are given in Theorem~\ref{thm:quotient-dynamics}.
\end{proposition}

\begin{proof}
Gauge invariance of every predicate sends an admitted state to another
admitted state.  The preservation witnesses make $F$ a self-map of
$\mathcal S_{\mathsf G}$, and the covariance equation remains true after this
restriction.  Hence gauge-equivalent inputs have gauge-equivalent outputs,
which defines the quotient update.  The relational preservation implication
removes every transition from an admitted state to a nonadmitted state.  The
stochastic preservation equality puts total probability one on admitted
outputs.  Equivariance is unchanged by either restriction.
\end{proof}

The three notions have different types.  A Gauss intertwiner is a
representation-theoretic vertex tensor.  A Gauss-state witness proves that a
complete field state satisfies a local constraint.  A Gauss-preservation
witness proves that a law carries such states to such states.  For example, a
finite field model may declare a gauge-invariant charge-balance function
$q_v:\mathcal S\to A$ into a finite Abelian tally group $A$, and take
$\mathsf{Gauss}_v(x)$ to mean $q_v(x)=0$.  In
Chapter~\ref{ch:5} the incidence operator $\dvg$ will be defined and the specialization
will be $q_v=\dvg E-\rho$, where $\rho$ is the charge of that chapter, not a
representation.  In a quantum representation model, the analogous
admissibility condition can be presented by projection to the invariant
subspace.  These are related realizations of local gauge neutrality, not one
untyped object.

\subsection{Why class labels of non-Abelian defects do not multiply}

A face defect may be labeled by the conjugacy class of its holonomy, that is,
by its image under the class map $\kappa:G\to\operatorname{Conj}(G)$ to the
set of conjugacy classes.  To fuse two defects, however, one must first choose
paths transporting their holonomy representatives to a common basepoint and
then multiply the transported representatives.  Class labels alone do not
retain this information \cite{PatternGeometry}.  The order of two factors does
not affect the class of their product, since $xy$ and $yx=x^{-1}(xy)x$ are
conjugate; from three defects on, the class can depend on their cyclic order.

\begin{example}[Why class labels do not multiply]
\label{ex:class-labels-do-not-multiply}
Let $A_5$ be the group of even permutations of five symbols.  Its fifteen
double transpositions form one conjugacy class, conventionally called $2A$.
Take $a=(12)(34)$ and $c=(12)(35)$.  Both therefore lie in $2A$.  Choosing
representatives $(a,a)$ gives product $1$, whereas $(a,c)$ gives a
three-cycle.  The input pair of class labels $(2A,2A)$ does not determine an
output class.  Transported representatives at a common basepoint are
indispensable: exchanging the two factors of either pair does not change the
class of the product, but replacing one representative by a conjugate, as a
different transport path does, can.
\end{example}

\begin{definition}[Witnessed defect fusion]
\label{def:defect-fusion-witness}
Let $f_1,f_2$ be oriented defect faces with selected boundary basepoints
$b_1,b_2$, let $\operatorname{Conj}(G)$ be the set of conjugacy classes, and
put $H_i=F_U^{b_i}(f_i)$.  A fusion witness over a chosen
comparison vertex $v_*$ is the record
\[
 w=(b_1,b_2,v_*,p_1,p_2,\epsilon):
   \AFFusionWitness_U(f_1,f_2;v_*),
\]
where $p_i:b_i\leadsto v_*$ are oriented transport paths and
$\epsilon$ is one of the two orders on $\{1,2\}$.  It determines transported
representatives and an output label
\[
 \widetilde H_i
 =\Hol_U(p_i)H_i\Hol_U(p_i)^{-1},
 \qquad
 \operatorname{Fuse}_U(w)
 =\bigl[\widetilde H_{\epsilon(2)}
         \widetilde H_{\epsilon(1)}\bigr]
 \in\operatorname{Conj}(G).
\]
The rightmost factor is fused first.  The basepoints, paths, and order are
fields of the WITNESS; the conjugacy class is its derived gauge-invariant
readout and therefore has RENDERING role.
\end{definition}

The class $\operatorname{Fuse}_U(w)$ does not depend on $\epsilon$, since
$\widetilde H_1\widetilde H_2
 =\widetilde H_1(\widetilde H_2\widetilde H_1)\widetilde H_1^{-1}$; the order
is kept because the ordered product
$\widetilde H_{\epsilon(2)}\widetilde H_{\epsilon(1)}$ itself depends on it.
The class does depend on the paths.  Replacing $p_1$ by another path $p_1'$
from $b_1$ to $v_*$ conjugates $\widetilde H_1$ by the holonomy of the loop
$p_1'p_1^{-1}$ at $v_*$, and this can change the class of the product, as the
representatives $a$ and $c$ of Example~\ref{ex:class-labels-do-not-multiply}
show.

\begin{proposition}[Gauge invariance of witnessed fusion]
\label{prop:defect-fusion-covariance}
\status{PROVED-LOCAL}
Gauge transformation carries a fusion witness to a witness on the transformed
state with the same paths and order, and
\[
 \operatorname{Fuse}_{\alpha U}(\alpha_*w)
 =\operatorname{Fuse}_U(w).
\]
If $G$ is Abelian, the result is independent of the chosen basepoints, paths,
and order and reduces to multiplication of the two defect labels.  For a
general group, forgetting those witness fields is not licensed.
\end{proposition}

\begin{proof}
Write $h_v$ for the coordinate value of $\alpha$ at $v$.  Endpoint covariance
of path holonomy and conjugation covariance of $H_i$ give
\[
 \widetilde H_i\longmapsto
 h_{v_*}\widetilde H_i h_{v_*}^{-1}.
\]
Their ordered product is therefore conjugated by the same $h_{v_*}$, so its
conjugacy class is unchanged.  In an Abelian group conjugation is trivial and
the two factors commute, proving all independence claims.  The converse, that
class labels alone determine the fused class for every pair only when $G$ is
Abelian, is Theorem~\ref{thm:centrality-package}\textup{(iv)}.
\end{proof}

When a dynamics joins defects, this record becomes the evidence field
of an application witness.  The application must still supply its geometric
match, guard, affected support, state update, and trace event.  Naming a fusion
observable does not by itself define that SCRIPT.

\section{The affine semidirect profile}
\label{sec:affine-textbook}

The preceding construction used one internal group.  Several gauge channels
may sometimes form a direct product, but they need not do so.  The affine
profile is the basic counterexample: a frame rotation acts on the meaning of
a translation, so the two channels must be coupled by a semidirect product.
The affine group and its covariance theorem are taken from
\cite{PatternGeometry}.  That source writes the affine group $\mathcal A$;
here $\mathcal A$ denotes the coefficient profile of
Definition~\ref{def:coefficient-profile-part1}, so the affine group is
written $\AffG$ and its translation group $\Lambda$.  In this section
$\Lambda$ is a group of translations, not the channel index of a coefficient
profile.

Let $\Lambda$ be a finite Abelian translation group, written additively;
let $F$ be a finite frame group acting by
$\rho_F:F\to\Aut(\Lambda)$; and let $G$ be a finite internal group.  Put
$K=F\times G$ and let $\rho:K\to\Aut(\Lambda)$ factor through $F$.

\begin{definition}[Affine semidirect gauge group]
The \emph{affine gauge group} is
\begin{equation}
\AffG=\Lambda\rtimes_\rho K,
 \qquad
 (t,k)(t',k')=(t+\rho(k)t',kk').
 \label{eq:affine-group}
\end{equation}
The first component records translation, while $k=(f,g)$ records frame and
internal data.  The action $\rho(k)$ in the multiplication rule says that the
first frame determines how the second translation is read.
\end{definition}

An $\AffG$-connection is no new kind of connection: it is the construction
of Section~\ref{sec:connections-local} with structure group $\AffG$.
Its loop holonomy jointly records failure of translational closure, frame
return, and internal-pose return.

\begin{example}[A finite translation--frame coupling]
\label{ex:affine-triangle}
Take $\Lambda=C_3$ additively, $F=C_2=\{0,1\}$, and let the nontrivial frame
element act by inversion, $\rho(1)t=-t$.  With trivial internal group,
$\AffG=C_3\rtimes C_2$ is the six-element symmetry group of a triangle.
Let $a=(1,0)$ be a unit translation and $b=(0,1)$ a frame reflection.  Then
\[
 bab^{-1}=(-1,0)=(2,0).
\]
Changing the frame has reversed the translation label.  Recording an
independent pair ``translation $1$, frame reflection'' would miss this
covariance and therefore would not define a stable raw defect label.
\end{example}

The following theorem is proved in \cite{PatternGeometry}; we prove it again
in the present notation.

\begin{theorem}[Affine covariance]
\label{thm:affine-covariance}
\status{PROVED-SOURCE+PROVED-LOCAL}
Conjugation in \eqref{eq:affine-group} is
\begin{equation}
(t,k)(t',k')(t,k)^{-1}
 =\bigl(t+\rho(k)t'-\rho(kk'k^{-1})t,\;kk'k^{-1}\bigr).
 \label{eq:affine-conjugation}
\end{equation}
Consequently the projected $K$-curvature transforms by conjugation, while the
translation component transforms affinely.  There is no general
gauge-invariant reading of the raw translation component as an independent
channel.  When $\rho$ is trivial, the channels decouple.
\end{theorem}

\begin{proof}
The inverse of $(t,k)$ is $(-\rho(k^{-1})t,k^{-1})$.  Hence
\[
\begin{aligned}
 (t,k)(t',k')(t,k)^{-1}
 &=(t+\rho(k)t',kk')(-\rho(k^{-1})t,k^{-1})\\
 &=\bigl(t+\rho(k)t'-\rho(kk'k^{-1})t,kk'k^{-1}\bigr).
\end{aligned}
\]
The second component is ordinary conjugation in $K$; the first mixes the old
translation $t'$, the gauge translation $t$, and the projected element $k'$.
For trivial $\rho$ this mixing becomes $t+t'-t=t'$.
\end{proof}

The correct invariant is the conjugacy class of the full
$\AffG$-holonomy.  A more refined description can be given after a
representative $k'$ of the projected $K$-class has been fixed.  Conjugation by
a pure translation changes $t'$ by an element of
$(1-\rho(k'))\Lambda$, meaning the image subgroup
\[
 (1-\rho(k'))\Lambda
 =\{s-\rho(k')s:s\in\Lambda\}.
\]
Thus its residual class lies in the
\emph{coinvariant quotient}
\[
 \Lambda_{k'}\defeq
 \Lambda/(1-\rho(k'))\Lambda.
\]
The \emph{centralizer}
$C_K(k')=\{k\in K:kk'=k'k\}$ still acts on this quotient, so the final datum is
a $C_K(k')$-orbit in $\Lambda_{k'}$.  Both the quotient and the residual
action are required; neither was visible in an independent-channel tuple.

\begin{example}[Coinvariants in $C_3\rtimes C_2$]
Return to the six-element triangle group of
Example~\ref{ex:affine-triangle}.  Over the projected
identity $k'=0$, the subgroup
$(1-\rho(0))C_3$ is $0$, so the coinvariant quotient is $C_3$.  The centralizer
$C_2$ acts by inversion and leaves two residual orbits,
$\{0\}$ and $\{1,2\}$.  Over the projected reflection $k'=1$,
$1-\rho(1)$ is multiplication by $2$.  This is a bijection of $C_3$, so the
image subgroup is all of $C_3$ and the quotient has one element.  Consequently
all three affine elements whose frame component is a reflection lie in one
conjugacy class.  The quotient calculation recovers the familiar conjugacy
classes of triangle rotations and reflections without treating translation
as an independent label.
\end{example}

\section{Gauge-covariant classical laws and profiles}
\label{sec:complete-gauge-profile}

\subsection{Coefficient profiles in intrinsic form}

Definition~\ref{def:coefficient-profile-part1} specifies how typed fields are
anchored to gauge sites.  A general gauge profile must import that interface,
not hide every nonconnection variable in an undifferentiated auxiliary factor
\cite{ArtC,PatternGeometry}.  What is added here is the intrinsic
associated-fibre form of the field type and the reduct that produces it.

\begin{definition}[Intrinsic field type and coefficient realization]
\label{def:coefficient-bridge}
Let
\[
 Q_X=(S_X,L_X,s,t,\overline{(-)})
\]
be a chosen transport skeleton of $X$ in the sense of
Definition~\ref{def:transport-skeleton-general}, and let
\[
 \mathcal A=
 \bigl(Q_X,G,(k_\lambda,A_\lambda,S_\lambda,
                    a_\lambda,\rho_\lambda)_{\lambda\in\Lambda}\bigr)
 :\AFCoefficientProfile(Q_X,G)
\]
be a coefficient profile on it (Definition~\ref{def:coefficient-profile-part1}):
each channel $\lambda$ of the finite index $\Lambda$ has cell dimension
$k_\lambda$, support $S_\lambda\subseteq C_{k_\lambda}(X)$, value object
$A_\lambda$, gauge-frame anchor $a_\lambda:S_\lambda\to S_X$, and declared
action $\rho_\lambda:G\to\Aut(A_\lambda)$.  Given a localization
$\mathcal T=(T_v)_v$, its intrinsic field type is
\[
 \AFField_{\mathcal T}(\mathcal A)
 =\AFDepFun_{\lambda:\Lambda}
   \AFDepFun_{x: S_\lambda}
   \bigl(T_{a_\lambda(x)}\times_G A_\lambda\bigr).
\]
The intrinsic gauge action on each factor is
$[t,b]\mapsto[\alpha_{a_\lambda(x)}(t),b]$ and acts componentwise on the
product.  This is the action meant by ``gauge-equivariant'' below.
Let $\mathcal M$ be any declared gauge-stable matter-state type and let
$\mathcal S_{\rm aux}$ be any declared auxiliary state type with a specified
gauge action.  A \emph{coefficient realization} of those two factors is a
gauge-equivariant reduct
\[
 r_{\mathcal A}:\mathcal M\times\mathcal S_{\rm aux}
       \longrightarrow\AFField_{\mathcal T}(\mathcal A).
\]
Its fibres state exactly which distinctions are forgotten by the coefficient
view.  When $r_{\mathcal A}$ is a bijection onto a declared admitted field
subtype, it is instead a coordinate-free PRESENTATION of all nonconnection
field state.
\end{definition}

\begin{definition}[Coefficient interpretation witness]
\label{def:coefficient-intwit}
Put $Z_{\mathcal A}=\mathcal M\times\mathcal S_{\rm aux}$.  For a typed map
$r:Z_{\mathcal A}\to\AFField_{\mathcal T}(\mathcal A)$, let
$\mathsf{CoeffLaw}_{\mathcal A}(r)$ contain its gauge-equivariance proof and
the kernel-pair receipt
\[
 z\sim_r z'\quad\Longleftrightarrow\quad r(z)=r(z').
\]
The typed codomain of $r$ certifies that it lands in the intrinsic field
type, and the kernel pair $\sim_r$ records exactly the distinctions that the
coefficient view forgets; both are determined by $r$.  Define
\[
 \AFCoeffIntWit_{\mathcal A}
   (Z_{\mathcal A},\AFField_{\mathcal T}(\mathcal A))
 \defeq
 \AFDepPair_{r:Z_{\mathcal A}\to\AFField_{\mathcal T}(\mathcal A)}
 \mathsf{CoeffLaw}_{\mathcal A}(r)
 =\AFIntWit_{\mathcal C^{\mathrm{coeff}}_{\mathcal A}}
   (Z_{\mathcal A},\AFField_{\mathcal T}(\mathcal A)).
\]
The coefficient profile $\mathcal A$ fixes the species of endpoint and law;
the actual reduct $r$ is the realizer field of an inhabitant and is not hidden
inside the profile index.
\end{definition}

\begin{proposition}[Coefficient-realization bridge]
\label{prop:coefficient-intwit-bridge}
First projection gives a canonical equivalence between coefficient
interpretation witnesses and the gauge-equivariant coefficient realizations
of Definition~\ref{def:coefficient-bridge}, with the target type and
forgotten fibres retained.  Write
$r_{\mathcal A}\defeq\operatorname{real}(\kappa_{\mathcal A})$ for this
projection from an inhabitant $\kappa_{\mathcal A}$.
\end{proposition}

\begin{proof}
A coefficient witness projects to its typed equivariant map.  Conversely,
such a map supplies the equivariance field, its codomain supplies target
agreement, and its equality fibres supply the uniquely determined kernel-pair
receipt.  These operations are inverse on the map and all extensional law
fields.
\end{proof}

This chapter and the next use the full skeleton,
\[
 Q_X=\mathsf T(X),\qquad S_X=V(X),\qquad L_X=E^{\rm or}(X),
\]
which is boundary-complete for every declared face of $X$.  Smaller transport
skeletons, and transfer of fields between different canvases, belong to a
typed geometric transition receipt, not to the fixed-canvas profile of
Definition~\ref{def:gauge-profile}.  In particular, scale transfer is neither
assumed nor discarded; it is additional data on an epoch or refinement arrow.

\begin{proposition}[Coefficient-interface compatibility]
\label{prop:coefficient-compatibility}
\status{PROVED-LOCAL}
Let $\kappa_{\mathcal A}$ be a coefficient interpretation witness and write
$r_{\mathcal A}=\operatorname{real}(\kappa_{\mathcal A})$.  Choose a local
trivialization of $\mathcal T$.  Then every intrinsic field in
$\AFField_{\mathcal T}(\mathcal A)$ has unique coordinates
$\phi_\lambda:S_\lambda\to A_\lambda$, and passive or active change acts by
\[
 (h\cdot\phi)_\lambda(x)
 =\rho_\lambda(h_{a_\lambda(x)})\phi_\lambda(x).
\]
Thus $(\mathcal A,\kappa_{\mathcal A})$, through its $r_{\mathcal A}$
projection, is exactly the coefficient interface of
Definition~\ref{def:coefficient-profile-part1} on the fixed canvas $X$.  In
particular, vertex
matter representations give the specialization
$k_v=0$, $S_v=\{v\}$, $a_v(v)=v$, and $A_v=M_v$.
\end{proposition}

\begin{proof}
For each associated fibre, a chosen $t_{a_\lambda(x)}$ writes its unique member
as $[t_{a_\lambda(x)},\phi_\lambda(x)]$.  The coordinate-change calculation is
the same one as Equation~\eqref{eq:matter-gauge-action}.  Taking the indicated
products proves the formula channel by channel.  The vertex specialization is
obtained by substituting the displayed support and anchor data.
\end{proof}

\subsection{Connecting the gauge state to the process language}

Let $\mathcal C\subseteq\Conn_G(X;(T_v)_v)$ be an admitted connection state
type stable under the intrinsic action of $\mathscr G_X$.  Let $\mathcal M$ be
a stable type of matter configurations in declared associated fibres, and let
$\mathcal S_{\rm aux}$ be an auxiliary state factor.  Every auxiliary
factor has a specified $\mathscr G_X$-action, which may be trivial.  The total
state type and action are
\begin{equation}
\mathcal S
 =\mathcal C\times\mathcal M\times\mathcal S_{\rm aux},
 \qquad
 \alpha\cdot(\mathcal U,\Psi,z)
 =(\alpha\cdot\mathcal U,\alpha\cdot\Psi,\alpha\cdot z).
 \label{eq:total-gauge-state}
\end{equation}
A finite palette, in the sense of Section~\ref{sec:matter-textbook}, makes a
representation-valued matter model finite-state in the sense of
Definition~\ref{def:gauge-profile}; a full nonzero complex vector space does
not.

Definition~\ref{def:controlled-law-local} defined a controlled law specification
$\mathfrak L=(\mathcal S,\mathcal I,\Theta,\lambda,\chi,\mathcal P)$ and kept
its scheduler policy separate from the specification.  State-level covariance
is necessary, but the retained process data require a finer definition.

\begin{definition}[Located application witness]
\label{def:located-gauge-application}
Let $\mathsf{Addr}(\mathcal S)$ be the finite type of addressable canvas and
field components of a fixed gauge state.  In the local presentation used by
this chapter, an inhabitant
$w\in\mathsf{App}_a(x)$ is the record
\[
 w=(p,\ell,R_w,W_w,A_w,g_w,r_w,q_w,c_w,y_w),
\]
where $p$ is a phase-$a$ rule name; $\ell$ is its located geometric match;
$R_w,W_w\subseteq\mathsf{Addr}(\mathcal S)$ are finite read and write
footprints; $A_w\supseteq R_w\cup W_w$ is its declared affected support;
$g_w$ proves the guard at $(x,\ell)$; $r_w=x|_{R_w}$ is the read receipt;
$q_w$ is the proposed assignment on $W_w$; and $c_w$ proves that committing
that proposal gives a legal state $y_w$ agreeing with $x$ outside $W_w$.
The semantics is $\operatorname{apply}_{a,x}(w)=y_w$, and
$\operatorname{foot}(w)=(R_w,W_w,A_w)$.
\end{definition}

This record specializes the application witnesses of
Definition~\ref{def:process-profile} to gauge states; it is not a second process formalism.
Figure~\ref{fig:c2-kick} shows its fields for the application of
Example~\ref{ex:complete-c2-profile}.  A batch
contains such witnesses together with its compatibility evidence.  Gauge
transport of an application preserves $p$, the underlying spatial location,
and the address sets, transports the match $\ell$ and
$g_w,r_w,q_w,c_w$ in their dependent fibres, and sends $y_w$ to
$\alpha y_w$.  A purely geometric match may transport by the identity; a
decorated match may do so only through its declared gauge-safe recognition
witness.  Thus a law is local when the finite footprint
$\operatorname{foot}(w)$ of each of its applications is exposed in the record.

\begin{definition}[Gauge action on a process profile]
\label{def:gauge-process-action}
Let $\mathcal P=(A,\mathsf{App},\operatorname{apply},
\mathsf{Batch},\operatorname{ApplyBatch})$ be a process profile over
$\mathcal S$.  The control alphabet and control tapestry are taken to be
gauge neutral.  A \emph{gauge action on $\mathcal P$} assigns, for every
$\alpha\in\mathscr G_X$, phase $a$, and state $x$, bijections
\[
 \alpha_*:\mathsf{App}_a(x)\longrightarrow
 \mathsf{App}_a(\alpha\cdot x),
 \qquad
 \alpha_*:\mathsf{Batch}_a(x)\longrightarrow
 \mathsf{Batch}_a(\alpha\cdot x).
\]
The application maps preserve rule and underlying location indices and the
three spatial footprints of
Definition~\ref{def:located-gauge-application}; they transport the located
match, guard, read, proposal, and commit receipts in their dependent fibres.  The
batch maps preserve the empty batch, order, and compatibility witnesses.
On both application and batch fibres the assignments obey
\[
 (1)_{*,x}=\id,
 \qquad
 (\alpha\beta)_{*,x}
 =\alpha_{*,\beta x}\circ\beta_{*,x},
\]
and they satisfy
\begin{align}
\operatorname{apply}_{a,\alpha x}(\alpha_*w)
   &=\alpha\cdot\operatorname{apply}_{a,x}(w), \label{eq:app-naturality}\\
 \operatorname{ApplyBatch}_{a,\alpha x}(\alpha_*B)
   &=\alpha\cdot\operatorname{ApplyBatch}_{a,x}(B).
 \label{eq:batch-naturality}
\end{align}
Thus gauge transformation transports not only the output state but also the
witness that licensed the update.
\end{definition}

\begin{definition}[Classical semantic modes]
\label{def:three-realizer-modes}
Fix a process profile with the action above.
\begin{enumerate}[label=\textup{(\roman*)},leftmargin=2.5em]
\item A deterministic policy $\sigma_a^{\rm det}$ induces
\[
 F_a(x)=\operatorname{ApplyBatch}_{a,x}
          \bigl(\sigma_a^{\rm det}(x)\bigr).
\]
\item A nondeterministic policy $\sigma_a^{\rm nd}(x)$, a nonempty subset of
$\mathsf{Batch}_a(x)$, induces the relation
\[
 (x,y)\in R_a
 \quad\Longleftrightarrow\quad
 \exists B\in\sigma_a^{\rm nd}(x)\ \
 \operatorname{ApplyBatch}_{a,x}(B)=y.
\]
\item A finitely supported stochastic policy
$\sigma_a^{\rm stoch}(x)\in
\mathsf{Dist}(\mathsf{Batch}_a(x))$ induces, on a finite state type, the
Markov kernel
\begin{equation}
K_a(x,y)=
 \sum_{\substack{B\in\mathsf{Batch}_a(x)\\
       \operatorname{ApplyBatch}_{a,x}(B)=y}}
       \sigma_a^{\rm stoch}(x)(B).
 \label{eq:batch-pushforward-kernel}
\end{equation}
\end{enumerate}
The \emph{classical semantic mode} records which of these three objects is
retained.  They are not three names for one untyped realizer.
\end{definition}

The three modes above exhaust the classical semantics used in Part~I.
Phase-weighted alternatives are not a fourth kind of classical policy.  They
are introduced only after the quantization bridge, in
Definition~\ref{def:coherent-realizer} of Chapter~\ref{ch:8}.

\begin{definition}[Process-level gauge covariance]
\label{def:process-gauge-covariance}
A policy is \emph{process-gauge-covariant} when, for all $\alpha,a,x$,
\begin{align*}
\sigma_a^{\rm det}(\alpha x)
   &=\alpha_*\sigma_a^{\rm det}(x),\\
 \sigma_a^{\rm nd}(\alpha x)
   &=\{\alpha_*B:B\in\sigma_a^{\rm nd}(x)\},\\
 \sigma_a^{\rm stoch}(\alpha x)(\alpha_*B)
   &=\sigma_a^{\rm stoch}(x)(B),
\end{align*}
using the line appropriate to its semantic mode.  If the realizer emits a
trace event $\operatorname{emit}_{a,x}(B)$, the trace encoder must additionally
carry a gauge action and satisfy
\[
 \operatorname{emit}_{a,\alpha x}(\alpha_*B)
 =\alpha_*\operatorname{emit}_{a,x}(B).
\]
A \emph{gauge-covariant realized law} consists of the process action, a
process-gauge-covariant policy, an equivariant trace encoder when traces are
retained, and the locality, totality, invariant-preservation, and
Gauss-preservation witnesses required by its mode and declared invariant
profile.  Batch-independence and realizer-independence are additional claims;
they belong to the run-contract slots below and are not consequences of gauge
covariance.
\end{definition}

\begin{proposition}[Process covariance implies extensional covariance]
\label{prop:process-covariance}
\status{PROVED-LOCAL}
A process-gauge-covariant deterministic policy induces
\begin{equation}
F_a(\alpha\cdot x)=\alpha\cdot F_a(x).
 \label{eq:law-gauge-covariance}
\end{equation}
A process-gauge-covariant nondeterministic policy induces
\[
 (x,y)\in R_a
 \quad\Longleftrightarrow\quad
 (\alpha\cdot x,\alpha\cdot y)\in R_a,
\]
and a process-gauge-covariant stochastic policy induces
\[
 K_a(\alpha\cdot x,\alpha\cdot y)=K_a(x,y).
\]
The converses need not hold if application witnesses, policies, or traces are
retained as structural data.
\end{proposition}

\begin{proof}
For the deterministic case, policy covariance followed by
Equation~\eqref{eq:batch-naturality} gives the displayed equation.  In the
relational case, a selected witness batch $B$ from $x$ is carried bijectively
to $\alpha_*B$ from $\alpha x$, whose output is $\alpha y$; applying
$\alpha^{-1}$ proves the converse implication.  In the stochastic case,
Equation~\eqref{eq:batch-pushforward-kernel} is unchanged after the bijective
reindexing $B\mapsto\alpha_*B$.

For failure of the converse, let two distinct application witnesses have the
same identity state update and let a scheduler choose between them by
inspecting a raw link coordinate.  The extensional update is the
gauge-covariant identity, while the retained choice and its trace can change
under gauge transformation.
\end{proof}

\begin{example}[A nontrivial gauge-covariant holonomy kick]
Fix a loop $\gamma$ based at $v$ and a representation $\rho:G\to\operatorname{GL}(W)$.
Let the state contain a connection and a matter vector at $v$.  Define a
one-step law that leaves the connection fixed and updates
\[
 \psi_v\longmapsto\rho(\Hol_U(\gamma))\psi_v,
\]
leaving all other matter coordinates fixed.  Under a gauge transformation,
Theorem~\ref{thm:holonomy-descent} gives
\[
 \rho(h_v\Hol_U(\gamma)h_v^{-1})\rho(h_v)\psi_v
 =\rho(h_v)\rho(\Hol_U(\gamma))\psi_v.
\]
Thus the update satisfies Equation~\eqref{eq:law-gauge-covariance}.  It is
nontrivial whenever the loop holonomy acts nontrivially, and it is reversible
because $\rho(\Hol_U(\gamma))$ is invertible.  Intrinsically, this is the
endomorphism of the associated fibre induced by the loop transport
$\mathcal U_\gamma:T_v\to T_v$.  A complete realized version must additionally
give the application, batch, policy, and trace actions of
Definition~\ref{def:process-gauge-covariance}.
\end{example}

\begin{definition}[Run contract]
\label{def:run-contract}
For a controlled realized law, a term
\[
 \mathfrak R:\AFRunContract(\mathfrak L,\mathsf{Real})
\]
contains a declared initial boundary
$\mathcal S_0\subseteq\mathcal S$; a mode-correct run assignment; its
$\AFTraceTapestry$ event type and encoder;
and the resulting solution-tapestry type synchronizing control, states,
events, and invariant witnesses.  In deterministic mode the run assignment
is the trajectory function $x_0\mapsto(x_0,\ldots,x_n)$; in relational mode
it is the nonempty type of admitted trajectories; and in stochastic mode it
is the finitely supported path distribution induced along the finite reading.
It also contains four explicit claim slots
\[
 \mathsf{CadFair},\quad \mathsf{AppFair},\quad
 \mathsf{BatchInd},\quad \mathsf{RealizerInd}.
\]
Each slot is either $\AFNoClaim$ or an inhabitant of the corresponding witness
type.  The first two certify cadence and application fairness; the third
certifies independence from the linear presentation of every compatible
batch; the fourth certifies confluence or observational equivalence across the
declared class of realizers.  A fairness witness must name the cyclic or
infinite extension on which recurrence is quantified; a finite reading by
itself proves neither kind of fairness.  Every independence witness also
states its class of batches, realizers, initial states, and observations.
Absence of a claim is therefore data, not a silently assumed theorem.
\end{definition}

Definition~\ref{def:run-contract} admits order-sensitive realized processes.
The deterministic \emph{physical realizer} of \cite{ToyElectrodynamics} is
stronger: it requires phasewise locality, selected-step totality,
batch-presentation independence, and invariant preservation.  In the present
profile such a model must inhabit $\mathsf{BatchInd}$ as well as the
locality, totality, and preservation fields.  Recording $\AFNoClaim$ widens
the class of admitted laws, and a law whose $\mathsf{BatchInd}$ slot is
$\AFNoClaim$ is not covered by the well-posedness theorem of
\cite{ToyElectrodynamics}.

The class of observations in $\mathsf{RealizerInd}$ must be stated.  In the
coherent theory, equality at every fixed physical background need not give
equality on a coherent joint input (Theorem~\ref{plaw:joint-channel-gluing}
in Chapter~\ref{ch:17}), and a restricted instrument library can leave
different states indistinguishable
(Theorem~\ref{plaw:finite-observation-history} in Chapter~\ref{ch:15}).
These criteria carry their own quantum hypotheses and do not follow from
locality or totality.

\begin{proposition}[Gauge covariance lifts to realized runs]
\label{prop:gauge-run-covariance}
\status{PROVED-LOCAL}
For a process-gauge-covariant realizer with equivariant trace encoder and a
finite control reading, gauge action lifts to its mode-correct run object:
\begin{enumerate}[label=\textup{(\roman*)},leftmargin=2.5em]
\item in deterministic mode it sends the run from $x_0$ to the run from
$\alpha\cdot x_0$ column by column;
\item in relational mode it gives a bijection between the admitted trajectory
and trace types from those two initial states;
\item in stochastic mode it pushes the path distribution from $x_0$ to the
path distribution from $\alpha\cdot x_0$, preserving the probability of each
gauge-related cylinder event.
\end{enumerate}
The configuration row is transformed pointwise.  In every mode
the event row uses its declared event action and every gauge-invariant
observable row is unchanged.  Hence a classical solution tapestry descends to
the initial-state gauge orbit, while its event-bearing tapestry need only
transform equivariantly.  This conclusion does not supply any fairness,
batch-independence, or realizer-independence witness absent from
$\mathfrak R$.
\end{proposition}

\begin{proof}
Induct over the chosen finite reading.  At one column, policy covariance maps
the selected batch, selected-batch set, or batch distribution according to the
mode; Equation~\eqref{eq:batch-naturality} maps its committed state, and trace
equivariance maps its emitted event.  In deterministic mode this gives the
next column.  In relational mode the bijection on selected batches gives a
bijection on one-step extensions, with inverse induced by $\alpha^{-1}$.  In
stochastic mode the same bijective reindexing preserves every one-step
conditional weight.  Iteration proves the trajectory statements; multiplying
the finitely many conditional weights proves the path-distribution statement.
Gauge invariance makes each observable row equal before and after transport.
The final warning follows because the induction uses none of the four claim
slots of Definition~\ref{def:run-contract}.
\end{proof}

Gauge covariance does not replace the other process obligations.  Conversely,
locality, totality, and invariant preservation do not imply gauge covariance.
An energy, action, or harmony score used by the scheduler must be a
gauge-invariant function of the admitted configuration, computed by the
policy with values in a declared formal type (the codomain criterion of
Section~\ref{sec:holonomy-textbook}); a displayed reading of the same
function is a RENDERING and cannot be consumed.  The score may rank already
legal applications, but the harmony firewall of Section~\ref{sec:harmony-firewall} forbids the ranking
from replacing the application witness or repairing an ill-typed update.  A
score may occur inside a guard only when that guard is independently declared
and its gauge covariance and other admissibility obligations are proved; it is
then part of the witness type, not a post hoc exception made by the scheduler.

\subsection{Assembling the finite classical gauge profile}

Before packaging the fields, one optional decoration must be typed.  A
\emph{pointed channel datum} for $G$ is either absent or consists of groups
$N,Q$ and maps
\[
 1\longrightarrow N\xrightarrow{i}G\xrightarrow{q}Q\longrightarrow1
\]
for which $i$ is injective, $q$ is surjective, and
$i(N)=\ker q$.  It may additionally retain a homomorphism
$s:Q\to G$ with $q\circ s=\id_Q$.  Section~\ref{sec:intrinsic-strata} will
classify these data as direct-split, genuinely semidirect-split, or nonsplit.

\begin{definition}[Finite-base and finite-state classical gauge profiles]
\label{def:gauge-profile}
A \emph{finite-base classical gauge profile} is the dependent record
\[
 \mathfrak G=(\mathfrak K,\mathfrak M,\mathfrak A,\mathfrak D;
                    \mathfrak P_G)
\]
with the following fields.
\begin{enumerate}[label=\textup{(\roman*)},leftmargin=2.5em]
\item The \emph{kinematic layer}
\[
 \mathfrak K=\bigl(X,Q_X=\mathsf T(X),G,
          \mathcal T=(T_v)_{v\in V(X)},
          \mathscr G_X,\mathcal C,\mathfrak E_{\bullet}\bigr)
\]
contains a finite oriented cellular canvas $X:\AFOrientedCellForm$, its full
transport skeleton, a finite group, a localization receipt
$\mathcal T:\AFLocalization_G(X)$, the intrinsic gauge group
$\mathscr G_X=\prod_v\Aut_G(T_v)$, a $\mathscr G_X$-stable connection type
$\mathcal C\subseteq\Conn_G(X;\mathcal T)$, and an optional pointed channel
datum $\mathfrak E_{\bullet}$.

\item The optional \emph{matter layer}
\[
 \mathfrak M=
 \bigl(\boldsymbol\rho,(E_v)_v,\mathcal M,
              (\mathcal I_v,\mathcal I_v^{\vee})_v\bigr),
 \qquad
 \boldsymbol\rho=
 \bigl((\rho_e)_{e\in E^+(X)},(\mu_v)_{v\in V(X)}\bigr),
\]
contains the declared edge and vertex representations, the associated fibres,
a $\mathscr G_X$-stable admitted matter-state type, and the invariant tensor
and contraction types.  A pure gauge model takes singleton matter state and
trivial carriers.  A finite-state representation model declares the finite
palette whose associated configurations form $\mathcal M$.

\item The \emph{coefficient layer}
\[
 \mathfrak A=(\mathcal S_{\rm aux},\mathcal A,\kappa_{\mathcal A})
\]
contains an auxiliary state type with its declared gauge action, a coefficient
profile on $Q_X$ with internal group exactly $G$, and
the coefficient interpretation witness
\[
 \kappa_{\mathcal A}:
 \AFCoeffIntWit_{\mathcal A}
  (\mathcal M\times\mathcal S_{\rm aux},
   \AFField_{\mathcal T}(\mathcal A))
\]
of Definition~\ref{def:coefficient-intwit}.  Its projection
$r_{\mathcal A}=\operatorname{real}(\kappa_{\mathcal A})$ is the coefficient
reduct of Proposition~\ref{prop:coefficient-intwit-bridge}, used below.  The
empty channel family is allowed.
Thus auxiliary state is never called a coefficient field without a declared
support, value object, anchor, action, and reduct map.

\item The \emph{observation and dynamics layer}
\[
 \mathfrak D=
 \bigl(\mathcal S,\mathsf{Rec},\mathsf{Obs},
       \mathsf{Gauss},\nu,
       \mathfrak L,\mathsf{Real},\mathfrak R\bigr)
\]
contains the required identity
\[
 \mathcal S=\mathcal C\times\mathcal M\times\mathcal S_{\rm aux}.
\]
The observable field is a typed family with a finite declared index type,
$\mathsf{Obs}=((Y_i,\omega_i))_{i\in I_{\rm obs}}$ with
\[
 \omega_i:\AFGaugeObsIntWit(\mathcal S,Y_i),
 \qquad O_i\defeq\operatorname{real}(\omega_i):\mathcal S\to Y_i.
\]
The recognition field is the finite family
\[
 \mathsf{Rec}=((\AFRecognitionProfile[j],\eta_j))_{j\in I_{\rm rec}},\qquad
 \eta_j:\AFGaugeSafeRecIntWit_{\AFRecognitionProfile[j],\mathcal S}.
\]
Definition~\ref{def:gauge-safe-recognition} gives the receipt
$w_j:\AFGaugeSafe(\AFRecognitionProfile[j],\mathcal S)$.  Either family may be empty.  The Gauss
field is a constraint profile on $\mathcal S$.  The field
$\nu\in\{\mathsf{det},\mathsf{rel},\mathsf{stoch}\}$ is the semantic mode of
Definition~\ref{def:three-realizer-modes}.  The controlled
law $\mathfrak L$ is required to have
configuration field exactly $\mathcal S$, and $\mathsf{Real}$ is a realizer
of mode $\nu$ carrying the process action, policy, trace encoder, and all
covariance and preservation witnesses required by that mode.  Finally,
$\mathfrak R:\AFRunContract(\mathfrak L,\mathsf{Real})$ exposes the trajectory,
trace, solution-tapestry, and optional fairness or independence claims.
\end{enumerate}
The final field is a typed source witness
$\mathfrak P_G:\AFGroupSource(G)$ from
Definition~\ref{def:group-source-receipt}.  It is retained in the construction
history but not used to compute the gauge action.

The profile is a \emph{finite-state classical gauge profile} when its total state type
$\mathcal S$ is finite.  Throughout Part~I, the unqualified phrase
\emph{finite classical gauge profile} means finite-base: it asserts finiteness of the
canvas and internal group, not of a vector-valued state space.  Every theorem
that counts states, forms a finite Markov sum, or uses finite-state dynamics
states the additional finite-state hypothesis explicitly.
\end{definition}

\begin{proposition}[Classical-profile compatibility projections]
\label{prop:gauge-profile-intwit-compatibility}
Projecting $\kappa_{\mathcal A}$ to $r_{\mathcal A}$, each $\omega_i$ to
$O_i$, and each $\eta_j$ to its receipt $w_j$ recovers an equivariant
coefficient reduct, a gauge-invariant observable family, and a gauge-safe
recognition family.  Conversely, such a reduct and such families package
canonically into the three interpretation-witness fields.  These operations
are mutually inverse on extensional data, preserve empty families, and leave
$\mathfrak K$, $\mathfrak M$, $\mathsf{Gauss}$, $\nu$, $\mathfrak L$,
$\mathsf{Real}$, $\mathfrak R$, and $\mathfrak P_G$ unchanged.
\end{proposition}

\begin{proof}
Apply Propositions~\ref{prop:coefficient-intwit-bridge}
and~\ref{prop:gauge-observable-intwit-bridge} componentwise and use the
definitional equality in Definition~\ref{def:gauge-safe-recognition}.  Finite
dependent products preserve these inverse maps, including the empty product.
All unlisted fields are copied literally.
\end{proof}

\begin{table}[!tb]
\centering
\footnotesize
\renewcommand{\arraystretch}{1.15}
\begin{tabularx}{\textwidth}{@{}>{\raggedright\arraybackslash}p{0.15\textwidth}
  >{\raggedright\arraybackslash}p{0.22\textwidth}
  >{\raggedright\arraybackslash}p{0.28\textwidth}
  >{\raggedright\arraybackslash}X@{}}
\toprule
\textbf{Field} & \textbf{From Chapter~\ref{ch:1}} &
\textbf{Added in Chapters~\ref{ch:2}--\ref{ch:3}} &
\textbf{In Example~\ref{ex:complete-c2-profile}}\\
\midrule
$X,Q_X=\mathsf T(X)$ & cellular canvas, orientation, boundary words & full
site/link transport specialization & the one-face square with $v_0,\ldots,v_3$,
$e_0,\ldots,e_3$, and $f$\\
$G,\mathfrak P_G,\mathcal T,\mathscr G_X$ & profiled automorphism group and
one marked-pose torsor & typed group source, independent local torsors, and
active fibre automorphisms & $C_2=\{1,z\}$; $\mathsf{Declared}(d_{C_2})$; a
regular torsor at each vertex; $\mathscr G_X\cong C_2^4$\\
$\mathcal C,\mathfrak E_{\bullet}$ & state-type vocabulary and declared
carried data & gauge-alignment witnesses, admitted connection states, and an
optional pointed channel & all $16$ connections; no pointed channel\\
$\boldsymbol\rho,\mathcal M,\mathcal I_v,\mathcal I_v^{\vee}$ &
representation and intertwiner vocabulary & associated matter fibres, finite
palettes, transport, invariant tensors, and contractions & the sign
representation at $v_0$ with palette $\{\pm1\}$, trivial elsewhere;
$\mathcal I_{v_0}=0$\\
$\mathcal A,\kappa_{\mathcal A}$ & coefficient profiles as parameter data,
not variable state & supports, anchors, values, actions, intrinsic field
type, and coefficient interpretation witness with reduct projection & a
matter channel and a charge channel at $v_0$; a bijective reduct\\
$\mathcal S_{\rm aux},\mathcal S$ & typed products and declared actions & the
total diagonal gauge action and its state-type equality & $q\in C_2$ with
trivial action; $16\cdot2\cdot2=64$ states\\
$\mathsf{Rec},\mathsf{Gauss},\mathsf{Obs}$ & recognition profiles, invariant
predicates, observables, and checkpoint witnesses & gauge-safe-recognition
and gauge-observable interpretation families, local constraints, and their
realizer projections & a profile reading incidence, curvature, and charge;
$q=F^{v_0}_{\mathcal U}(f)$ at $v_0$; $O_{\rm flux}$ and $O_{\rm charge}$\\
$\nu,\mathfrak L,\mathsf{Real},\mathfrak R$ & process, cadence, policy,
trace, and invariants & mode-correct process action, gauge covariance, runs,
fairness claim slots, and Gauss preservation & $\mathsf{det}$; the one phase
$\mathsf{KICK}$ and the singleton batch $(w_x)$; a one-step run whose two
fairness slots are $\AFNoClaim$\\
\bottomrule
\end{tabularx}
\caption{The fields of Definition~\ref{def:gauge-profile}: what
Chapter~\ref{ch:1} provides for each, what Chapters~\ref{ch:2}
and~\ref{ch:3} add, and the value of each in the complete $C_2$ profile of
Example~\ref{ex:complete-c2-profile}.}
\label{tab:gauge-profile-fields}
\end{table}

Table~\ref{tab:gauge-profile-fields} lists the fields of the record, with
what each chapter contributes and their values in the complete $C_2$ example.
The dependencies in this record are part of its type.  In particular,
$\mathfrak L$ cannot silently use a state type different from
Equation~\eqref{eq:total-gauge-state}, a stochastic realizer cannot inhabit a
deterministic mode, and a pointed channel cannot be inferred from the abstract
group merely because the group has a normal subgroup.  Likewise, a raw
recognition profile cannot enter $\mathsf{Rec}$ without its gauge-safety
witness, and an auxiliary value does not enter $\mathcal A$ without an anchor,
action, and explicit coefficient reduct.

\begin{definition}[Artform gauge-frame enrichment]
\label{def:artform-gauge-frame}
For a finite-base classical gauge profile $\mathfrak G$, the expression
\[
 \widehat{\mathfrak G}=\AFGaugeFrame(\mathfrak G):\AFArtformFrame
\]
is the frame enrichment whose constructed carrier is the already formed canvas
$X$, whose carried structure is the dependent profile record, and whose
construction history is
\[
 \AFHistory(\widehat{\mathfrak G})
 =\bigl(\AFHistory(X),
        \mathsf{GaugeProfileNode}(\mathfrak G)\bigr).
\]
The dependent profile is installed PRESENTATION structure on $X$.  Its
dynamical state type is
\[
 \AFState(\widehat{\mathfrak G})=\mathcal S.
\]
Thus $\mathcal S$ is the typed configuration state consumed by located
applications; its inhabitants are not additional role indices.  Its WITNESS
fields include group-source and torsor receipts, alignment, covariance,
constraint, recognition, compatibility, preservation, and run certificates
when applicable.
Its SCRIPT fields are active gauge transformations, admitted process
applications, realized classical updates when they have SCRIPT type, and the
forgetful reduct maps defined below.  PRESENTATION data
also include installed local trivializations, coordinate conventions, state
types, and the realizer policy once installed; link and coefficient tables are
inhabitants of those state types.  The realizer projections of
$\mathsf{Obs}$ and the
descended occurrence statistics licensed by $\mathsf{Rec}$ are terminal
readouts with RENDERING role.  A diagram or table displaying them is likewise
a RENDERING; nothing in the formal profile may read either result back as
input.
\end{definition}

Figure~\ref{fig:gauge-five-roles} sorts these fields by role.

\begin{figure}[htbp]
\centering
\begin{tikzpicture}[
  box/.style={draw=black!50,rounded corners=2pt,align=center,
              minimum width=40mm,minimum height=10mm,font=\scriptsize}]
\node[box,fill=tealrod!12] (form) at (0,0)
  {FORM / FRAME\\$X,\widehat{\mathfrak G}$};
\node[box,fill=purplerod!10] (pres) at (0,1.55)
  {PRESENTATION\\profile, $\mathcal S$, coordinates, policy};
\node[box,fill=goldrod!13] (wit) at (-3.45,3.1)
  {WITNESS\\source, alignment, guards};
\node[box,fill=paperbg] (script) at (3.45,3.1)
  {SCRIPT\\gauge/update/reduct maps};
\node[box,fill=greenrod!12] (render) at (0,4.65)
  {RENDERING\\class, character, tally, display};
\draw[artarrow] (form)--(pres);
\draw[artarrow] (pres)--(wit);
\draw[artarrow] (pres)--(script);
\draw[artarrow] (wit)--(render);
\draw[artarrow] (script)--(render);
\draw[designarrow,rounded corners=5pt]
  (wit.west)--(-5.55,3.1)--(-5.55,0)--
  node[below,font=\scriptsize]{registered gate} (form.west);
\end{tikzpicture}
\caption{The five Artform roles specialized to a gauge frame.  Solid arrows
show selected ordinary dependencies and the dashed return is gate-mediated.
Configuration variables such as $\mathcal U,\Psi,z$ are not role nodes: they
inhabit state types carried by the PRESENTATION.  A policy moves from
WITNESS to installed PRESENTATION structure, reduct maps remain SCRIPTS, and
all readouts are terminal RENDERINGS.  No arrow leaves RENDERING.}
\label{fig:gauge-five-roles}
\end{figure}

Define the base reduct
\[
 \operatorname{ForgetGauge}(\widehat{\mathfrak G})=X
\]
and the configuration-level coefficient reduct $\operatorname{Chr}$, which
sends a state of $\widehat{\mathfrak G}$ to
\[
 \operatorname{Chr}(\mathcal U,\Psi,z)=
 (X,\mathcal A,r_{\mathcal A}(\Psi,z),\mathcal U).
\]
We call its target $\mathfrak C_{\mathfrak G}$ the \emph{coefficient
presentation} of $\mathfrak G$ and write
$\operatorname{Chr}:\widehat{\mathfrak G}\to\mathfrak C_{\mathfrak G}$.  A
value of $\operatorname{Chr}$ is the intrinsic fixed-canvas version of a
$\AFChromoForm$ state (Definition~\ref{def:coefficient-profile-part1}); a
local trivialization presents its field and connection coordinates.  The
reduct forgets information exactly when $r_{\mathcal A}$ has nontrivial
fibres.

\begin{theorem}[Gauge-frame formation and base section]
\label{thm:gauge-frame-conservativity}
\status{PROVED-LOCAL}
Every record satisfying Definition~\ref{def:gauge-profile} forms the Artform
frame $\widehat{\mathfrak G}$.  The base reduct preserves literally the cell
inventory, grading, incidence, orientation, boundary words, and construction
history of $X$.  Moreover every finite oriented canvas has a trivial gauge
enrichment $\operatorname{TrivGauge}(X)$ for which
\[
 \operatorname{ForgetGauge}(\operatorname{TrivGauge}(X))=X.
\]
Consequently gauge-frame formation adds no new core FORM constructor: the
canvas layer is a retract of the gauge layer at the level of objects.
\end{theorem}

\begin{proof}
The first statement is dependent-record introduction using the fields already
listed.  The forgetful map is projection to the first field, so all listed
canvas data and $\AFHistory(X)$ are unchanged.  For the section, take the
trivial group, its declared source receipt, the unique one-point torsor at
every vertex, the singleton connection, matter, and auxiliary state types, the
empty coefficient and recognition families, true Gauss predicates, and a
one-phase identity process with an empty invariant profile, its unique
deterministic realizer, and a run contract making no fairness or independence
claim.  Every required equation has its unique witness.  Projection returns
$X$.  No field of the record is a cell, an incidence, or a boundary word of
$X$, so no core FORM constructor was used to build it.
\end{proof}

\subsection{Syntactic conservativity of the gauge increment}

A section on objects is not yet a conservativity theorem.  Whether the
enrichment can prove a new equation \emph{between paths of the canvas
fragment} is a question about syntax, and the foundational calculus answers
it with a declared finite check rather than an appeal to the shape of the
construction \cite{ArtC}.  In this subsection $\mathcal M$ and $\mathcal M'$
denote equational theories and $\rho$ denotes a shadow map; they are
unrelated to the matter-state type $\mathcal M$ and to representations.  The
criterion for a language extension
$\mathcal M\subseteq\mathcal M'$ asks for a typed graph map
\[
 \rho:\mathcal G_{\mathcal M'}\longrightarrow
      \mathsf{Path}(\mathcal G_{\mathcal M}),
\]
the \emph{old shadow}, subject to two conditions: \textup{(R1)} $\rho$ is the
identity on old vertices and old generators, and \textup{(R2)} the induced
functor $\rho^\sharp$ sends both sides of every equation of the increment to
paths that are already provably equal in the old theory.  Such a $\rho$
exhibits the old syntactic category as a retract of the new one and therefore
makes the extension syntactically conservative.  The criterion is sufficient,
not necessary; and it consults no state algebra, a point we return to below.
Definition~\ref{def:gauge-increment} lists the increment so that the
criterion can be checked.

\begin{definition}[The gauge increment $\Delta\mathcal M_{\rm gauge}$]
\label{def:gauge-increment}
Let $\mathcal M_{\rm can}$ be the canvas fragment built from the constructors
of Chapter~\ref{ch:1}: its vertices are
the oriented cellular-canvas expressions, and its generators are the
FORM-producing constructors $\AFAttachCell,\AFRefine,\AFQuotient$ together
with the forgetful reducts between them.  The increment
$\Delta\mathcal M_{\rm gauge}$ adjoins, over each canvas expression $X$, the
following data and nothing else.

\emph{Vertices.}  For every finite-base classical gauge profile $\mathfrak G$ over $X$,
the frame $\widehat{\mathfrak G}$; its coordinate presentation
$\widehat{\mathfrak G}_t$ for each local trivialization $t$; its extensional
structural reduct $\widehat{\mathfrak G}^{\rm ext}$; and the coefficient
presentation $\mathfrak C_{\mathfrak G}$, the target of $\operatorname{Chr}$.
Here $\widehat{\mathfrak G}_t$ and $\widehat{\mathfrak G}^{\rm ext}$ are the
targets of the maps $R_{\rm coord}^t$ and $R_{\rm ext}$ of the reduct ladder
(Definition~\ref{def:gauge-reduct-ladder}).

\emph{Generators.}  The enrichment $\operatorname{TrivGauge}_X$; the base
reducts $\pi_A:A\to X$ from each new vertex $A$, with
$\pi_{\widehat{\mathfrak G}}=\operatorname{ForgetGauge}$; the ladder maps
$R_{\rm ext}$, $R_{\rm coord}^t$ and its inverse $\Phi_t^{-1}$, where $\Phi_t$
is the coordinate presentation of Theorem~\ref{thm:localization}, and the
coefficient reduct $\operatorname{Chr}$; one active gauge script
$\gamma_\alpha:\widehat{\mathfrak G}\to\widehat{\mathfrak G}$ for each
$\alpha\in\mathscr G_X$; and, in deterministic mode, one realized update
$\AFStep_a:\widehat{\mathfrak G}\to\widehat{\mathfrak G}$ for each phase $a$.
In relational and stochastic mode the realized law is not an endomorphism
script edge and contributes no such generator; every clause below is
unaffected.

\emph{Equations.}  Writing $A,B$ for new vertices and $e:A\to B$ for a new
generator between them,
\[
\begin{aligned}
 &\text{(V1)}&\pi_B\circ e&\equiv\pi_A,&
 &\text{(V2)}&\operatorname{ForgetGauge}\circ\operatorname{TrivGauge}_X
   &\equiv 1_X,\\
 &\text{(V3)}&\gamma_1&\equiv 1_{\widehat{\mathfrak G}},&
 &\text{(V4)}&\gamma_\alpha\circ\gamma_\beta&\equiv\gamma_{\alpha\beta},\\
 &\text{(V5)}&\AFStep_a\circ\gamma_\alpha&\equiv\gamma_\alpha\circ\AFStep_a,&
 &\text{(V6)}&\Phi_t^{-1}\circ R_{\rm coord}^t
   &\equiv 1_{\widehat{\mathfrak G}},\\
 &\text{(V7)}&R_{\rm coord}^t\circ\Phi_t^{-1}
   &\equiv 1_{\widehat{\mathfrak G}_t}.&&&&
\end{aligned}
\]
Clause (V1) is the content of
Proposition~\ref{prop:localization-formation} and of the rule that only an
operation on the state of an already fixed frame is a SCRIPT: no new generator
moves the canvas.  Clauses (V3)--(V4) say that
Equation~\eqref{eq:intrinsic-gauge-action} is a group action, (V5) is
Equation~\eqref{eq:law-gauge-covariance} in deterministic mode; (V6)--(V7) are
Theorem~\ref{thm:localization}\textup{(i)}.

No observable, occurrence tally, character, or orbit count is a generator of
$\Delta\mathcal M_{\rm gauge}$.  That is not an omission but the readback
firewall: a RENDERING is terminal, so no composite may leave it, and a
terminal readout can never occur in a parallel pair of old paths.
\end{definition}

\begin{proposition}[The gauge increment is syntactically conservative]
\label{prop:gauge-retraction}
\status{PROVED-LOCAL}
Send every new vertex over $X$ to $X$, and every new generator to the empty
path at its shadow:
\[
 \rho_0(\widehat{\mathfrak G})
 =\rho_0(\widehat{\mathfrak G}_t)
 =\rho_0(\widehat{\mathfrak G}^{\rm ext})
 =\rho_0(\mathfrak C_{\mathfrak G})
 =X,
 \qquad
 \rho_1(e)=1_X .
\]
Then $\rho$ is a reduct retraction of
$\mathcal M_{\rm can}\subseteq\mathcal M_{\rm can}+\Delta\mathcal M_{\rm gauge}$.
Consequently the gauge increment is syntactically conservative over the canvas
fragment: two parallel paths of $\mathcal M_{\rm can}$ become equal in the
extended theory only if they were already equal in $\mathcal M_{\rm can}$.
\end{proposition}

\begin{proof}
Typing first.  Every generator listed in
Definition~\ref{def:gauge-increment} has both endpoints over the same canvas
expression $X$: the ladder maps, the coefficient reduct, and the gauge and
update scripts run between new vertices over $X$, while
$\operatorname{TrivGauge}_X$ and the base reducts $\pi_A$ have $X$ at one end
and a new vertex over $X$ at the other.  After shadowing, each therefore runs
from $X$ to $X$, which is the domain and codomain of the empty path $1_X$.
Condition~\textup{(R1)} holds by construction, since $\rho$ was declared only
on the new vertices and generators.

For~\textup{(R2)}, apply $\rho^\sharp$ to each equation of the increment.  A
composite of new generators maps to a composite of empty paths, hence to
$1_X$.  Both sides of (V1)--(V7) are such composites, or are an identity
already, so in every case both images are literally $1_X$.  Literal equality
of paths is in particular provable equality in $E_{\mathcal M_{\rm can}}$.  The
inherited old equations are mapped to themselves by~\textup{(R1)} and are old
provable equalities.  Thus $\rho$ is a reduct retraction, and the foundational
retraction theorem gives the syntactic conclusion \cite{ArtC}.
\end{proof}

\begin{remark}[The retraction is not state-sound]
\label{rem:not-state-sound}
The foundational criterion consults no state algebra, and this retraction
would fail the stricter \emph{state-sound} condition, which additionally
demands a bijection between the legal states of a new expression and those of
its old shadow.  Here $\AFState(\widehat{\mathfrak G})$ is the configuration
state $\mathcal S$, while the bare canvas $X$ carries no gauge state at all;
the one-square $C_2$ profile of
Example~\ref{ex:complete-c2-profile} has $64$ total states over a canvas with
none.  So the gauge increment is a genuine addition of state that is
nevertheless syntactically conservative on the canvas fragment.  That is the exact
content of the claim, and it is weaker than a re-presentation of old states:
Proposition~\ref{prop:gauge-retraction} licenses no transfer of a canvas-level
theorem to gauge states, and no state-level completeness.  The criterion is
also only sufficient, so a failure to produce $\rho$ for some later increment
would be evidence about the method, not about that increment.
\end{remark}

Table~\ref{tab:gauge-enrichment-stages} lists what each stage of the
construction adds.  On a fixed canvas $\AFGaugeFrame$ is an enrichment and
adds no carrier constructor; only a change of canvas uses $\AFAttachCell$ or
$\AFRefine$.  The conservativity claim is Proposition~\ref{prop:gauge-retraction},
and Remark~\ref{rem:not-state-sound} states its limits.  $\AFGaugeFrame$ is a
named enrichment because it exposes an important domain reading, not because
gauge theory needs a new universal carrier constructor.

\begin{table}[!tbp]
\centering
\footnotesize
\renewcommand{\arraystretch}{1.12}
\begin{tabularx}{\textwidth}{@{}>{\bfseries\raggedright\arraybackslash}p{.14\textwidth}
  >{\raggedright\arraybackslash}X>{\raggedright\arraybackslash}X
  >{\raggedright\arraybackslash}X>{\raggedright\arraybackslash}X@{}}
\toprule
Stage & Added structure & Added FORM constructors & Witnesses / scripts & Derived public operators\\
\midrule
source and localization & $\mathfrak P_G$, $\mathcal T$, tagged fibre support & none on $X$; dependent frame enrichment only & source, torsor, and localization receipts & $\mathscr G_X$, local frame types\\
connection and coefficients & $\mathcal C$, $\mathcal A$, matter and auxiliary actions & none & gauge alignments, invariant tensors, coefficient interpretation witnesses & holonomy, curvature, associated transport\\
recognition and observation & $\mathsf{Rec}$, $\mathsf{Obs}$, Gauss profile & none & gauge-safe-recognition, gauge-observable, and constraint witnesses & descended occurrences, tallies, orbit readings\\
process and runs & law, realizer, $\mathfrak R$ & none & applications, guards, footprints, compatibility, covariance, preservation & trajectories, trace and solution tapestries, quotient laws\\
canvas or scale transition & not a fixed-frame field & $\AFAttachCell$ or $\AFRefine$ when declared & geometric and field-transfer receipt & transported profile at the target epoch\\
\bottomrule
\end{tabularx}
\caption{What gauge enrichment adds at each stage of the construction.  Only
a change of canvas uses a FORM-producing constructor.}
\label{tab:gauge-enrichment-stages}
\end{table}

\subsection{A complete finite example}

\begin{figure}[htbp]
\centering
\begin{tikzpicture}[x=1cm,y=1cm,
  vtx/.style={circle,fill=platenavy,inner sep=1.5pt},
  rd/.style={draw=tealrod,line width=2.2pt,-{Latex[length=2.4mm]}},
  tok/.style={draw=orangerod,line width=1.2pt,fill=orangerod!18,rounded corners=2pt,
              font=\scriptsize,inner sep=2pt},
  aux/.style={draw=plategray,dashed,rounded corners=2pt,fill=white,font=\scriptsize,inner sep=2pt}]
\fill[tealrod!7] (0,0) rectangle (2.6,2.3);
\draw[rd] (0,0)--(2.6,0) node[midway,below,plate label]{$e_0$};
\draw[rd] (2.6,0)--(2.6,2.3) node[midway,right,plate label]{$e_1$};
\draw[rd] (2.6,2.3)--(0,2.3) node[midway,above,plate label]{$e_2$};
\draw[rd] (0,2.3)--(0,0) node[midway,left,plate label]{$e_3$};
\foreach \p in {(0,0),(2.6,0),(2.6,2.3),(0,2.3)}{\node[vtx] at \p {};}
\node[font=\scriptsize,anchor=south west] at (2.6,2.3) {$v_2$};
\node[font=\scriptsize,anchor=south east] at (0,2.3) {$v_3$};
\node[font=\scriptsize,anchor=north west] at (2.6,0) {$v_1$};
\node[font=\small] at (1.3,1.15) {$f$};
\node[tok,anchor=north east] (psi) at (-.12,-.12) {$\psi_{v_0}\in\{+1,-1\}$};
\node[font=\scriptsize,anchor=south east] at (-.05,.05) {$v_0$};
\node[aux,anchor=north] at (1.3,-.72) {charge $q\in C_2$};
\node[font=\scriptsize,anchor=west,align=left] at (4.0,2.25)
  {\textbf{rule} $\mathsf{KICK}$ at $(\partial_{v_0}f,v_0)$; guard true};
\node[font=\scriptsize,anchor=west,align=left,text=tealrod!80!black] at (4.0,1.6)
  {\textbf{read} $R_w$: the four boundary links\\and the matter value at $v_0$};
\node[font=\scriptsize,anchor=west,align=left,text=orangerod!85!black] at (4.0,.8)
  {\textbf{write} $W_w$: the matter value at $v_0$;\\proposal $\psi_{v_0}\mapsto\varepsilon(F^{v_0}_U(f))\,\psi_{v_0}$};
\node[font=\scriptsize,anchor=west,align=left,text=plategray] at (4.0,-.05)
  {\textbf{untouched}: the links and the charge $q$;\\so $\mathsf{Gauss}_{v_0}:\ q=F^{v_0}_{\mathcal U}(f)$ is preserved};
\end{tikzpicture}
\caption{The $\mathsf{KICK}$ application of Example~\ref{ex:complete-c2-profile}
on its square. Teal links form the read footprint; the orange token at
$v_0$ is read and written; the charge $q$, drawn dashed, belongs to the state
but not to the application's footprint. The proposal multiplies the sign
matter value by the character of the face curvature. Links and charge are
not written, so the Gauss predicate that compares them survives the commit.}
\label{fig:c2-kick}
\end{figure}

\begin{example}[A complete finite $C_2$ profile with a nontrivial law]
\label{ex:complete-c2-profile}
Take the one-face square with vertices $v_0,\ldots,v_3$, edges
$e_i:v_i\to v_{i+1}$ (indices modulo $4$) and face $f$ with boundary word
$(e_0e_1e_2e_3)$; Figure~\ref{fig:c2-kick} draws it with the data below.
Take $G=C_2=\{1,z\}$, the receipt $\mathfrak P_G=\mathsf{Declared}(d_{C_2})$,
the localization consisting of one regular torsor at every vertex, and all
torsor connections as $\mathcal C$; take the pointed-channel field to be
absent.  Choose a boundary vertex $v_0$ and the sign
representation $\varepsilon(z)=-1$ on $M_{v_0}=\mathbb C$.  Its finite stable
palette $P=\{+1,-1\}$ supplies a two-element associated matter fibre at
$v_0$; take $\mu_{v_0}=\varepsilon$ and all edge and other vertex
representations to be trivial.  All other matter-state factors are
singletons.  Let the auxiliary state be $q\in C_2$ with trivial gauge action.
The total state is
\[
 \mathcal S=\mathcal C\times
 (T_{v_0}\times_{C_2}P)\times C_2.
\]
Take two coefficient channels.  The matter channel has
$k_{\rm m}=0$, support $\{v_0\}$, value object $P$, anchor $v_0$, and sign
action.  The charge channel has the same support and anchor, value object
$C_2$, and trivial action.  The coefficient reduct sends $(\Psi,q)$ to these
two intrinsic associated-fibre values; it is a bijection here.  Together with
its equivariance and kernel-pair receipt it supplies
$\kappa_{\mathcal A}$, whose realizer projection is this same reduct.

Because $C_2$ is Abelian, the loop automorphism of $T_{v_0}$ has a canonical
value $F_{\mathcal U}^{v_0}(f)\in C_2$, independent of a local origin.  At
$v_0$ impose the nontrivial Gauss predicate
\[
 \mathsf{Gauss}_{v_0}(\mathcal U,\Psi,q)
 \quad\Longleftrightarrow\quad
 q=F_{\mathcal U}^{v_0}(f),
\]
and take the other vertex predicates to be true.  Both sides are gauge
invariant because $C_2$ is Abelian, and both are formal data of the state:
$q$ is a state component and $F_{\mathcal U}^{v_0}(f)$ a derived state.  The
predicate compares the auxiliary charge with the curvature value in $C_2$, not
with any reading of it.  This is what licenses the Gauss-preservation
obligation below to depend on the clause.

Let $\mathsf{Obs}$ contain two gauge-observable interpretation witnesses
whose realizer projections are the terminal readings
\[
 O_{\rm flux}=\kappa\circ F_{\mathcal U}^{v_0}(f),
 \qquad
 O_{\rm charge}=\kappa\circ q,
\]
into $\operatorname{Conj}(C_2)$, with their gauge-invariance proofs as law
certificates.  Here $\kappa$ is the class map of
Section~\ref{sec:matter-textbook}, not the witness $\kappa_{\mathcal A}$.  Let
$\mathsf{Rec}$ contain the recognition profile that sees
incidence, the curvature value, and the charge value but no raw link
coordinate, paired with its gauge-safe-recognition interpretation witness.
Both selected decorations are formal state data, so the profile may consume
them; neither $O_{\rm flux}$ nor $O_{\rm charge}$ is consumed anywhere.  The
identity position match is gauge-safe, so
Proposition~\ref{prop:gauge-recognition-descent} applies.

For the process data, take the one-element phase alphabet
$A=\{\mathsf{KICK}\}$, a one-tick control tapestry whose only column decodes
to $\mathsf{KICK}$, the reading $\lambda=(0)$, and the corresponding constant
decoder $\chi$.  Declare the semantic mode to be deterministic.  Let the
controlled-law invariant profile require the joint Gauss-state witness
$\AFDepFun_v\mathsf{Gauss}_v$ at the input and output checkpoints.  At every state let
$\mathsf{App}_{\mathsf{KICK}}(x)$ contain one witnessed application
$w_x$ with the following exposed fields:
\begin{itemize}[leftmargin=2.2em]
\item the rule is $\mathsf{KICK}$ and its location is the pair
      $(\partial_{v_0}f,v_0)$;
\item the guard is true; the read footprint consists of the four boundary
      connection components and the matter component at $v_0$; the write
      footprint is that matter component; and the affected support is their
      union;
\item the read receipt is the indicated restriction of $x$, and the proposal
      replaces $\Psi_{v_0}$ by the intrinsic associated-fibre transport
      $(\mathcal U_{\partial f})_*\Psi_{v_0}$ while fixing every other
      component;
\item the commit witness certifies that this proposal is in $\mathcal S$ and
      agrees with $x$ outside the write footprint.
\end{itemize}
In any local trivialization the same SCRIPT is
\[
 \operatorname{apply}(w_x):
 \psi_{v_0}\longmapsto
 \varepsilon(F_U^{v_0}(f))\psi_{v_0},
\]
with $U$ and $q$ fixed.  The batch type contains the required empty batch and
the singleton batch $(w_x)$; the deterministic policy selects $(w_x)$.
Compatibility and batch-presentation independence are immediate because no
batch contains two nonempty applications.  Gauge transformation sends every
field of $w_x$ to the corresponding field of $w_{\alpha x}$ while leaving its
spatial footprints fixed, so application and batch naturality hold.  The trace
event is
\[
 (0,\mathsf{KICK},w_x,\operatorname{foot}(w_x),x,
                  \operatorname{apply}(w_x));
\]
it retains the intrinsic loop transport rather than a preferred raw link
table.  Figure~\ref{fig:c2-seven-lane-run} reads this tick through the seven
lanes of the process reading of Section~\ref{sec:cadence-policy}.

\begin{figure}[htbp]
\centering
\small
\renewcommand{\arraystretch}{1.15}
\begin{tabularx}{.94\textwidth}{@{}>{\bfseries\raggedright\arraybackslash}p{.17\textwidth}X@{}}
\toprule
Clock & tick $0$ of the one-point reading $\lambda$\\
Control & decoded phase $\mathsf{KICK}$\\
Enabled witness & $w_x$ with true guard and located loop $\partial_{v_0}f$\\
Policy & select the singleton batch $(w_x)$\\
Proposal & multiply the $v_0$ matter coordinate by the intrinsic face transport\\
Atomic commit & replace only $\Psi_{v_0}$; preserve $\mathcal U$, $q$, and Gauss\\
Trace & record time, phase, witness, footprints, and the two committed states\\
\bottomrule
\end{tabularx}
\caption{The realized run of the complete $C_2$ profile of
Example~\ref{ex:complete-c2-profile}, read through the seven lanes of the
process reading of Section~\ref{sec:cadence-policy}.  Each row is a different
typed component even though all seven refer to the same tick; several
components may occupy the same one of the five Artform roles.}
\label{fig:c2-seven-lane-run}
\end{figure}

These data define the deterministic realizer: locality is the displayed
affected support, totality is supplied by the unique witnessed application at
every state, batch independence is the singleton argument, gauge covariance
is the preceding naturality square, and invariant and Gauss preservation hold
because the SCRIPT changes neither $\mathcal U$ nor $q$.

For the run contract take $\mathcal S_0=\mathcal S_{\mathsf G}$.  It assigns
the one-step trajectory, the displayed one-event trace tapestry, and their
solution tapestry.  Its cadence- and application-fairness slots are
$\AFNoClaim$, because a one-tick reading makes no recurrence assertion.  Its
batch-independence slot contains the preceding singleton proof, and its
realizer-independence slot is witnessed relative to the declared singleton
class containing this policy.

The law is reversible, is nonidentity on states with flux $z$, and preserves
the Gauss predicate because both $\mathcal U$ and $q$ remain fixed.  It is
nevertheless gauge-vertical: the constant gauge transformation $z$ fixes every
$C_2$ link and flips the sign matter coordinate, so the pre- and post-kick
states lie in one gauge orbit.  This example inhabits every field of
Definition~\ref{def:gauge-profile}, and Table~\ref{tab:gauge-profile-fields}
collects the values.  After all classifiers have been
introduced, Example~\ref{ex:complete-c2-dossier} will compute its intrinsic
dossier without using any additional profile data.
\end{example}

